\documentclass[aspectratio=169]{beamer}
\usetheme{Madrid}
\usecolortheme{beaver}
\usepackage[utf8]{inputenc}
\usepackage[T1]{fontenc}
\usepackage{tikz}
\usepackage{amssymb}
\usepackage{graphicx}
\setbeamertemplate{navigation symbols}{}
\setbeamercolor{title}{fg=white,bg=cyan!60!black}
\setbeamercolor{frametitle}{fg=white,bg=cyan!50!black}
\setbeamercolor{block title}{bg=cyan!15,fg=black}
\setbeamercolor{block body}{bg=cyan!5,fg=black}

\title{Block Arcade 1 — Sprites \& the Stage}
\subtitle{Every sprite lives at an (x, y) — coordinates in blocks}
\author{ByteForge --- Grade 7 --- Game Maker --- Week Overview}
\date{}

\begin{document}

\begin{frame}
    \titlepage
\end{frame}

\begin{frame}{This Week's Plan}
    \begin{itemize}
        \item Lesson 1: Meet Your Sprite
        \item Lesson 2: Two Sprites, Two Corners
        \item[$\bigstar$] BOSS --- Place Your Cast (\texttt{lessonBoss.tex})
    \end{itemize}
\end{frame}

\begin{frame}{Real-World Example}
    \Large
    Open Scratch's own stage and drag a sprite around — the x/y readout in the corner is literally the same two numbers this chapter is teaching. Minecraft goes further still: press F3 and you'll see your OWN character's live X, Y, Z coordinates on screen.
\end{frame}


\begin{frame}{Let Them Discover (1/2)}
    \Large
    On a printed 480x260 grid, mark where you'd place a sprite for "top-left corner", "dead center", and "bottom-right corner" before touching the platform — compare answers as a class.
\end{frame}


\begin{frame}{Let Them Discover (2/2)}
    \Large
    If you have access to Minecraft, open the F3 debug screen and walk in one direction. Which of X/Y/Z changes? Which stays the same? (Y is height in Minecraft — different from this chapter's 2D y, worth pointing out.)
\end{frame}

\begin{frame}{Files in This Folder}
    \small
    \begin{itemize}
        \item \texttt{lesson1.tex} --- Meet Your Sprite
        \item \texttt{lesson2.tex} --- Two Sprites, Two Corners
        \item \texttt{lessonBoss.tex} --- BOSS: Place Your Cast
    \end{itemize}
    \vspace{0.5em}
    \normalsize Each lesson file is its own short deck --- open the one for
    the specific day you're teaching, not this whole overview.
\end{frame}

\end{document}
